\documentclass[10.5pt,letterpaper]{article}

\usepackage[T1]{fontenc}
\usepackage{lmodern}
\usepackage[margin=0.75in]{geometry}
\usepackage{parskip}
\usepackage{enumitem}
\usepackage{xcolor}
\usepackage{titlesec}
\usepackage{fancyhdr}
\usepackage{microtype}

\setlength{\parskip}{4pt}
\setlist[itemize]{leftmargin=1.1em,topsep=1pt,itemsep=1pt,parsep=0pt}
\setlist[enumerate]{leftmargin=1.3em,topsep=1pt,itemsep=1pt,parsep=0pt}

\titleformat{\section}{\large\bfseries}{}{0pt}{}
\titlespacing*{\section}{0pt}{8pt}{2pt}
\titleformat{\subsection}{\normalsize\bfseries}{}{0pt}{}
\titlespacing*{\subsection}{0pt}{4pt}{1pt}

\pagestyle{fancy}
\fancyhf{}
\rhead{\small Microbiology LIMS POC --- v3.2}
\lhead{\small Process Reference}
\rfoot{\small \thepage}

\begin{document}

\begin{center}
{\Large\bfseries Microbiology LIMS POC --- Process at a Glance}\\[2pt]
{\small v3.2 (Conditioned Discovery) \quad\textbar\quad Reference to \texttt{microbio\_lims\_architecture.drawio} \quad\textbar\quad 2026-04-29}
\end{center}

\section{What the system does}

A subject-matter expert (SME) drops microbiology SOPs into S3. An agent pipeline reads them and produces two things for the dev team: \textbf{(A)} Jira-ready \emph{Epics} and \emph{Stories}, and \textbf{(B)} a per-culture \emph{configuration profile} in YAML (one each for urine, blood, target-pathogen).

The existing Cytology Jira backlog is used as a \textbf{teaching corpus} --- example (SOP excerpt $\rightarrow$ Story) pairs the agent learns the house style from. It is \emph{not} a registry to match against. The agent stops at Story. \emph{Tasks} (test IDs, library versions, DB columns) are dev-authored --- the agent has no codebase visibility to fabricate them defensibly.

\section{The pipeline --- 5 agents, run as a single batch over the SOP cohort}

\begin{enumerate}
  \item \textbf{Epic Extractor} reads the cohort and identifies capability domains, running in v3.2 \emph{conditioned mode} (matches drafts against the prior epic catalog before clustering anything novel; see \S v3.2 below). \emph{Out:} \texttt{\{id, title, discipline, themes[], theme\_catalog\_version, epic\_analog?\}}. Example title: \emph{Microbiology Specimen Receipt}.
  \item \textbf{Story Extractor} reads each SOP and emits Stories. Each Story declares one of \textbf{four shapes}, because real SOPs produce a mix --- a single rubric over-fits:
    \begin{itemize}
      \item \textbf{Capability} (a configurable feature; uses MUST/SHALL) --- e.g.\ \emph{Specimen Rejection Rules Engine}.
      \item \textbf{Workflow-stage-split} (same capability, decomposed by stage; siblings cross-linked) --- e.g.\ \emph{PHI Update Before Work Has Started / After Work Has Started / After Results Reported}.
      \item \textbf{Configuration-instance} (concrete values for one culture; this is data, not capability language) --- e.g.\ \emph{Urine centrifugation: speed 1500 g, duration 5 min, temp 4\,$^\circ$C}.
      \item \textbf{Cleanup / correction} (modify or remove an existing artifact) --- e.g.\ \emph{Remove Result Bucket option from Add/Modify QT code screens}.
    \end{itemize}
  \item \textbf{Validator --- gate \#1} verifies the declared shape, then applies a \emph{shape-specific} rubric (capability requires MUST/SHALL + testable AC; config-instance requires concrete typed values + verbatim citation; cleanup requires named target + before/after + regression risk; stage-split requires stage in title + sibling enumeration). Two revision rounds; on exhaustion, escalate to SME with the failed checks visible.
  \item \textbf{Cross-SOP Synthesis} runs only after gate \#1 has cleared the entire batch. It clusters concrete stories whose members come from \textbf{$\geq 2$ distinct SOPs} (two stories from the same SOP do \emph{not} qualify) and emits one capability-shaped story per cluster, with \texttt{cross\_links[]} pointing back at the contributing concrete stories. \textbf{Concrete stories are kept} --- the capability story is added as a sibling, not a replacement.
  \item \textbf{Validator --- gate \#2} (same agent as gate \#1) re-validates the synthesized stories against the capability rubric --- they are a \emph{new} artifact, not a derivative of validated input.
  \item \textbf{Dependency Resolver} --- final pass. Resolves story-to-story dependencies, stage-split sibling links, and capability-to-child links from synthesis. Outputs Jira-ready artifacts.
\end{enumerate}

\subsection{Story schema (what gets pushed to Jira)}

\noindent\texttt{\small \{\\
\hspace*{1em}id, epic\_id, title, description,\\
\hspace*{1em}shape $\in$ \{capability, workflow-stage-split, configuration-instance, cleanup\},\\
\hspace*{1em}acceptance\_criteria[\,], source\_citations[\,], dependencies[\,], cross\_links[\,],\\
\hspace*{1em}estimated\_complexity $\in$ \{S, M, L\}, source\_chunks[\,]\\
\}}

\section{The substrate --- how the agents see source material}

\textbf{Ingestion} (offline, fires when an SOP lands in S3): \emph{parse} (PDF / OCR for scans / DOCX / transcripts / Jira CSV) $\rightarrow$ \emph{chunk} (section-aware, $\sim$150 chunks per typical SOP) $\rightarrow$ \emph{enrich} (clinical NER for organisms / drugs / specimens, code expansion via UMLS, PHI detection \& redaction, theme tag from G1--G8) $\rightarrow$ \emph{embed} (both meaning-based vectors and keyword tokens for BM25) $\rightarrow$ \emph{store} (OpenSearch for the indexes, DynamoDB for chunk + doc metadata, S3 keeps the originals as the citation target).

\textbf{Retrieval} (real-time, per chat query): normalize $\rightarrow$ \emph{intent decode} into \texttt{\{action, themes, discipline, granularity, confidence\}} $\rightarrow$ expand with synonyms $\rightarrow$ \emph{multi-slot retrieve} (in-theme target / Cyto analogy / adjacent themes / Jira exemplars / regulatory --- each slot runs keyword + meaning search and fuses results) $\rightarrow$ rerank $\rightarrow$ assemble a role-labeled prompt (\textsc{Primary} / \textsc{Analogy} / \textsc{Adjacent} / \textsc{Exemplar} / \textsc{Reg}) $\rightarrow$ LLM generation with citations resolved to S3 doc + section + version.

Two thresholds gate the chat to prevent hallucination:
\begin{itemize}
  \item $\tau = 0.7$ \emph{intent confidence} --- below this, ask a clarifying question instead of retrieving.
  \item $\theta = 0.5$ \emph{retrieval relevance} --- if every slot falls below, refuse and offer the user three options (upload SOPs, draft from analogy with caveat, clarify scope).
\end{itemize}

\section{Worked example --- one query end-to-end}

\textbf{Query:} \emph{``Generate epics for specimen rejection in microbiology.''}

\textbf{Intent decode:} \texttt{\{action: generate, themes: [G1], discipline: Micro, granularity: epic, conf: 0.91\}} $\Rightarrow$ proceed.

\textbf{Retrieve $\rightarrow$ rerank:} 48 candidate chunks across 5 slots, trimmed to 20 by the reranker; low-score gate passes.

\textbf{Assemble:} 6{,}840-token prompt with the chunks role-labeled and the output schema attached.

\textbf{LLM output (structured JSON, abbreviated):}

\noindent\texttt{\small \{ epic\_id: epic\_draft\_3, title: "Specimen Rejection Workflow", theme: G1,\\
\hspace*{1em}citations: ["s3://.../SOP-MICRO-007 \S AC", "s3://.../SOP-CYTO-021 \S Procedure", "CLSI M22 \S 3.4"],\\
\hspace*{1em}stories: ["Validate transport temperature", "Flag specimens older than 4 h", \ldots] \}}

\textbf{Telemetry written:} query id, intent, slot scores, total chunks (20), tokens (6{,}840), latency ($\sim$4.3 s).

\section{What ships}

\begin{itemize}
  \item \textbf{Output A --- Jira artifacts.} Epics (with optional \texttt{cyto\_epic\_analog} annotation), per-SOP concrete Stories in all four shapes, synthesized capability Stories cross-linked to their children, and the dependency edges between them.
  \item \textbf{Output B --- Per-culture YAML profile.} One profile per culture in scope. Every value is typed (units / enums / nullable rules) and cites the SOP excerpt it came from. Inline references to the related Story's AC give traceability without duplicating values into the Story body.
\end{itemize}

\textbf{Example YAML excerpt} (\texttt{cultures/urine.yaml}):

\noindent\texttt{\small centrifugation:\\
\hspace*{1em}speed:    \{value: 1500, units: g\}    \# cite: SOP-MICRO-014 \S 3.2 L41--43\\
\hspace*{1em}duration: \{value: 5,\ \ \ units: min\}  \# cite: SOP-MICRO-014 \S 3.2 L44\\
\hspace*{1em}temp:     \{value: 4,\ \ \ units: degC\} \# cite: SOP-MICRO-014 \S 3.2 L45\\
related\_story\_ac: STORY-1042\#AC-2}

\section{v3.2 --- Conditioned Discovery (warm-start theme \& epic bootstrap)}

v3.1 generated epics from-scratch and treated themes G1--G8 as a closed Cyto-derived vocabulary. v3.2 adds a \emph{warm-start} step at both levels.

\textbf{(1) Theme Discovery agent (new, pre-flight).} Runs once per discipline (or when the G0 alarm trips). Four passes: \textbf{Classify} chunks against the prior catalog at $\tau_{\text{match}}$ ($0.65$) $\rightarrow$ \textbf{Cluster} only the residuals (size $\geq \varepsilon_{\text{novelty}}=3$) $\rightarrow$ SME \textbf{ratifies} novel candidates (inherited bypass) $\rightarrow$ \textbf{Re-tag} the corpus on catalog version bump. Emits \texttt{<discipline>\_v<N>}.

\textbf{(2) Epic Extractor in conditioned mode.} For each draft epic: classify against the prior epic catalog (\texttt{cyto\_epic\_v1}); if score $\geq \tau_{\text{match}}$, link \texttt{epic\_analog} and reuse the existing epic ID by construction; otherwise the residual is clustered across SOPs and SME ratifies novel ones. Emits \texttt{<discipline>\_epic\_v<N>}.

\textbf{Result.} Most themes and epics inherit their structure with explicit lineage; only genuinely-novel structure (with cluster evidence) goes to SME. The dev team reads the output and immediately sees what's reused vs.\ new.

\textbf{Worked example --- bootstrapping Histology from \texttt{cyto\_v1}:} 460/600 chunks classify into G1--G8 (residual $= 140$). Pass 2 surfaces 4 novel themes (H1 Tissue Processing, H2 Staining, H3 Block Archive, H4 Frozen Section); 22 chunks $\rightarrow$ G0 (3.7\%, below alarm). Pass 3 SME confirms all 4; flags G7 as redundant in histo (folded into H1) --- discarded. Pass 4 re-classifies the 12 G7-orphaned chunks (most land in H1). Final \texttt{histo\_v1}: 11 active themes (7 inherited + 4 novel; G7 in ANALOGY map only). 7/8 of the prior themes survived.

\textbf{ANALOGY map artifact (new).} Explicit cross-catalog link table replacing v3.1's implicit shared-theme assumption. Each link typed: \texttt{identical} / \texttt{partial} / \texttt{discarded\_in\_target} / \texttt{novel\_in\_target}. Drives the ANALOGY retrieval slot at query time.

\textbf{Micro POC migration.} Snapshot Cyto's themes as \texttt{cyto\_v1} and Cyto's epic backlog as \texttt{cyto\_epic\_v1}. Run conditioned Epic Extractor over the 3 in-scope SOPs to produce \texttt{micro\_epic\_v1} (most epics inherit; AST is the likely novel). Stories carry \texttt{theme\_catalog\_version: cyto\_v1} (Micro reuses the Cyto theme catalog wholesale; only the epic catalog forks). Emit a \texttt{cyto\_epic\_v1 $\leftrightarrow$ micro\_epic\_v1} ANALOGY map alongside.

\section{Cross-cutting (one line)}

KMS encryption at rest \& in transit; IAM least-privilege with VPC endpoints; DLQ for parse / OCR / embed failures; CloudWatch metrics + alarms; CloudTrail audit log on every API call. \emph{Generalization (v3.2):} when Histology arrives, Theme Discovery warm-starts \texttt{histo\_v1} from \texttt{cyto\_v1} pre-flight; the main pipeline then runs unchanged in shape.

\end{document}
